\section{A discriminant-dependent part of the real Hecke gate}
\label{sec:s5-real-zero-free}

This section proves a part of $(\mathcal R_B)$ using the specified
Dirichlet theorem $\mathsf A$, without the general Hecke hypothesis
$M$.  It does not prove $(\mathcal R_B)$ on the remaining central
interval.  The notation is that of MA-S5:
$L/\mathbb Q$ is the specified $S_5$ extension,
$E=L^{A_5}$, $F=L^{F_{20}}$, $H=EF$,
$C=L(s,\eta)$, and $B=\zeta_H/\zeta_F=C L_W$.

\paragraph{MA-RB-001 (an explicit endpoint exclusion).}
Put
\[
 \Delta_L=|\operatorname{disc}L|,\qquad
 \mathcal H_L=1+\tfrac12\log\Delta_L,\qquad
 \kappa_5=(\sqrt5-1)^2,\qquad
 \delta_L=\min\{\tfrac12,\kappa_5/\mathcal H_L\}.
 \tag{MA-RB.1}
\]
Under $\mathsf A$, the quadratic Hecke function $B$ has no real zero in
\[
              (0,\delta_L)\ \cup\ (1-\delta_L,1).
 \tag{MA-RB.2}
\]
More precisely, if $1/2<\sigma_0<1$ and $B(\sigma_0)=0$, then
\[
        1-\sigma_0\ge\kappa_5/\mathcal H_L.
 \tag{MA-RB.3}
\]
The endpoints in (MA-RB.2) are open.  No assertion about a possible
central zero, or about any remaining real point, is included.

\begin{proof}
\emph{The signed Artin orders at the particular point.}
Fix a real $\sigma_0\in(0,1)\setminus\{1/2\}$ and suppose that
$r=\operatorname{ord}_{\sigma_0}B\ge1$.
The theorem $\mathsf A$ gives $C(\sigma_0)\ne0$.
Consequently the all-prime identity $B=C L_W$ gives
\[
                 \operatorname{ord}_{\sigma_0}L_W=r.
 \tag{MA-RB.4}
\]
This is a pointwise deduction from a quotient of holomorphic
functions with nonzero denominator, and does not assume global
holomorphy of the unknown five-dimensional Artin factor.

The other five-dimensional factor is also holomorphic at this
particular real point:
\[
                   L_{W\eta}=\zeta_F/\zeta_{\mathbb Q}.
 \tag{MA-RB.5}
\]
Indeed $\zeta_F$ is holomorphic on $(0,1)$, and
$\zeta_{\mathbb Q}(\sigma)<0$ there.  For the latter sign, the
convergent alternating Dirichlet series is strictly positive
for real $\sigma>0$, while its factor
$1-2^{1-\sigma}$ is negative when $0<\sigma<1$.
Calegari's Theorem~1.1 supplies holomorphy of the irreducible
four- and six-dimensional factors $L_V,L_{V\eta},L_X$.
The regular-representation factorization, including every
ramified Euler factor, is
\[
 \zeta_L
   =\zeta_{\mathbb Q}\,C\,
       L_V^4L_{V\eta}^4L_W^5L_{W\eta}^5L_X^6.
 \tag{MA-RB.6}
\]
Taking signed orders at $\sigma_0$, the first two orders are
zero and every other order on the right is nonnegative.
It follows that
\[
                 \operatorname{ord}_{\sigma_0}\zeta_L
                         \ge5r\ge5.
 \tag{MA-RB.7}
\]
In particular no unproved holomorphy of either
five-dimensional factor has entered this inequality.
Equations (MA-RB.4)--(MA-RB.6) use the already checked
MT-S506 (S506.6), MT-S508 (S508.2), and the regular-character
identity in MT-S507 (S507.3).

\emph{Completion, signature, and the exact Hadamard identity.}
Because $L/\mathbb Q$ is Galois and its complex conjugation is the
nonidentity element of type $(12)(34)$, no embedding of $L$ is
real.  Its degree is $|S_5|=120$, so its signature is
$(r_1,r_2)=(0,60)$.  With
\[
 \Gamma_{\mathbb R}(s)=\pi^{-s/2}\Gamma(s/2),
 \qquad
 \Gamma_{\mathbb C}(s)=2(2\pi)^{-s}\Gamma(s),
\]
take the completed entire function
\[
 \xi_L(s)
   =s(s-1)\Delta_L^{s/2}\Gamma_{\mathbb C}(s)^{60}\zeta_L(s).
 \tag{MA-RB.8}
\]
The multiplicative constant in this normalization is immaterial.
The usual Dedekind continuation and functional equation give
$\xi_L(s)=\xi_L(1-s)$ and
$\overline{\xi_L(\overline s)}=\xi_L(s)$.
The simple pole and nonzero residue of $\zeta_L$ at $1$ imply
$\xi_L(1)\ne0$; the functional equation gives $\xi_L(0)\ne0$.
The zeros of $\xi_L$ are the nontrivial zeros, with
$0\le\operatorname{Re}\rho\le1$, including their multiplicities.
Only this containing strip is needed.

Here are the needed growth and normalization details of the
Hadamard formula.  On $\operatorname{Re}s\ge2$ the Euler product
bounds $\zeta_L(s)$, and the gamma estimate gives
$|\xi_L(s)|\le\exp(C|s|\log(|s|+3))$.
The functional equation gives the same bound on
$\operatorname{Re}s\le-1$.
On the remaining fixed strip, the unconditional polynomial
bound for the finite Dedekind function, with its pole removed,
and Stirling's formula give the same bound.
These are the fixed-strip inputs explicitly established in
NT-Z10 and recalled in MT-M18.
Thus $\xi_L$ is entire of order at most one.
The same local zero count gives
$\sum_{|\rho|>1}|\rho|^{-2}<\infty$.
Hadamard factorization therefore has the form
\[
 \xi_L(s)=\xi_L(0)e^{a_Ls}
        \prod_\rho(1-s/\rho)e^{s/\rho},
 \qquad
 \frac{\xi_L'}{\xi_L}(s)
   =a_L+\sum_\rho\left(\frac1{s-\rho}+\frac1\rho\right).
 \tag{MA-RB.9}
\]
The series in the second formula converges locally normally
off the zeros.  The constant
$a_L=\xi_L'(0)/\xi_L(0)$ is real.
Moreover
$S_L=\sum_\rho\operatorname{Re}(1/\rho)$ converges absolutely:
for large $|\rho|$ each term has absolute value at most
$|\rho|^{-2}$, by the containing strip.
Differentiating the functional equation at $0$ and $1$ gives
$\xi_L'(1)/\xi_L(1)=-a_L$.
Taking real parts of (MA-RB.9) at $1$ and using the
multiset symmetry $\rho\mapsto1-\rho$ gives
$-a_L=a_L+2S_L$.
Hence $a_L=-S_L$.  Taking real parts at a real $\sigma>1$
now proves, with its constant fixed rather than suppressed,
\[
 \frac{\xi_L'}{\xi_L}(\sigma)
             =\sum_\rho\operatorname{Re}\frac1{\sigma-\rho}.
 \tag{MA-RB.10}
\]
All summands are nonnegative.  The displayed real series
converges absolutely, again by the containing strip and the
zero count.  This proves the particular Hadamard identity
used here from the standard factorization theorem.

\emph{The archimedean estimate with all signs retained.}
Logarithmically differentiating (MA-RB.8) on the real
half-line $\sigma>1$ gives the exact formula
\[
 \sum_\rho\operatorname{Re}\frac1{\sigma-\rho}
 =\frac1{\sigma}+\frac1{\sigma-1}
       +\frac12\log\Delta_L
       +60\bigl(\psi(\sigma)-\log(2\pi)\bigr)
       +\frac{\zeta_L'}{\zeta_L}(\sigma),
 \tag{MA-RB.11}
\]
where $\psi=\Gamma'/\Gamma$.
The last term is negative, since the absolutely convergent
Dedekind Euler product gives
\[
 -\frac{\zeta_L'}{\zeta_L}(\sigma)
    =\sum_{\mathfrak p}\sum_{j\ge1}
          \frac{\log N\mathfrak p}{(N\mathfrak p)^{j\sigma}}>0.
 \tag{MA-RB.12}
\]
For completeness, $\psi(x)\le\log x$ for $x>0$ follows
directly from the gamma integral.  The density
$u^{x-1}e^{-u}/\Gamma(x)$ on $(0,\infty)$ has expectation
of $u$ equal to $x$, and expectation of $\log u$ equal
to $\psi(x)$; Jensen's inequality for the concave logarithm
gives the claim.  Differentiation of the gamma integral is
justified on compact positive $x$ intervals by integrable
majorants at zero and infinity.
Therefore, when $1<\sigma\le2$,
\[
                \psi(\sigma)-\log(2\pi)
                    \le\log\bigl(\sigma/(2\pi)\bigr)<0.
\]
Writing $\sigma=1+t$, $0<t\le1$, and dropping only
nonpositive terms in (MA-RB.11), we obtain
\[
        \sum_\rho\operatorname{Re}\frac1{1+t-\rho}
                     \le \frac1t+\mathcal H_L.
 \tag{MA-RB.13}
\]

\emph{Optimization and reflection.}
Suppose now that $1/2<\sigma_0<1$ is a zero of $B$, and put
$d=1-\sigma_0$.  Its at least five occurrences in the
zero divisor of $\xi_L$ contribute at least $5/(t+d)$
to (MA-RB.13).  Consequently
\[
                  \frac5{t+d}\le\frac1t+\mathcal H_L
                      \qquad(0<t\le1).
 \tag{MA-RB.14}
\]
Here $\Delta_L>1$ follows even from the displayed formulas:
at $\sigma=2$, (MA-RB.11)--(MA-RB.12) and the gamma bound give
$0\le3/2+\frac12\log\Delta_L-60\log\pi$, which excludes
$\Delta_L=1$.  Since the discriminant is an integer,
$\mathcal H_L\ge1+\frac12\log2>\sqrt5-1$.
Thus $t=(\sqrt5-1)/\mathcal H_L$ is permitted.
Substitution in (MA-RB.14) gives
$d\ge(\sqrt5-1)^2/\mathcal H_L$, proving (MA-RB.3).

To reflect the result, the quadratic Hecke completion of $B$
has root number $+1$, as the quotient of the completed
Dedekind equations for $H$ and $F$.
Its conductor factor is
$(|\operatorname{disc}H|/|\operatorname{disc}F|)^{s/2}$,
and its archimedean factor is
\[
                \Gamma_{\mathbb R}(s)^4
                    \Gamma_{\mathbb R}(s+1)^2.
 \tag{MA-RB.15}
\]
Indeed the settled permutation character
$\operatorname{Ind}_{F_{20}}^{S_5}\mathbf1=\mathbf1+W\eta$
has value two at complex conjugation, giving signature
$(2,2)$ for $F$; since $E$ is real, the signature of $H$
is $(4,4)$.  Taking the ratio of their gamma factors yields
(MA-RB.15), equivalently $\Gamma_{\mathbb R}(s)^2
\Gamma_{\mathbb C}(s)^2$.
All these factors are finite and nonzero on $(0,1)$.
The completed functional equation therefore pairs a real
zero $\sigma_0$ with $1-\sigma_0$, of the same order.
Reflecting the upper exclusion and using the definition
of $\delta_L$ proves (MA-RB.2).
\end{proof}

\paragraph{The exact remaining obstruction.}
The proof uses $\mathsf A$ only to make the quadratic Dirichlet
factor $C$ nonzero at the specified off-center point.
It supplies no nonvanishing assertion for the quadratic Hecke
function $B$ over the nonabelian sextic field $F$ in
$[\delta_L,1-\delta_L]\setminus\{1/2\}$.
The fact that its character is pulled back from $\eta$ does
not turn its $L$-function over $F$ into a Dirichlet
$L$-function over $\mathbb Q$: induction gives
$\operatorname{Ind}_{G_F}^{G_{\mathbb Q}}\eta_F=\eta\oplus W$,
which is precisely the unresolved nonabelian factorization.

The simple-real-zero descent theorem does not fill that interval.
Heilbronn (1973), Theorem~1, printed p.~870, says that a simple
real zero of the Dedekind zeta function of a normal extension
descends to the compositum of its quadratic subfields.
For $L/\mathbb Q$ this compositum is $E$, and $\mathsf A$
excludes such a zero of $\zeta_E$ off the center.
But a zero of $B$ in the present argument already has
multiplicity at least five in $\zeta_L$, by (MA-RB.7);
it is outside that theorem's simplicity hypothesis.

There is no contradiction from the signed order identities alone.
For example, at a hypothetical reflected pair of real points one
may assign
\[
 (\operatorname{ord}\zeta_{\mathbb Q},
   \operatorname{ord}C,\operatorname{ord}L_V,
   \operatorname{ord}L_{V\eta},\operatorname{ord}L_W,
   \operatorname{ord}L_{W\eta},\operatorname{ord}L_X)
                       =(0,0,0,0,1,0,0).
 \tag{MA-RB.16}
\]
Every induced-character order is then nonnegative, because
it is a nonnegative multiplicity times the sole nonzero
irreducible order.  The regular order is five and the
quadratic-subfield order is zero.  Assigning the same data
at the reflected point is compatible with functional-equation
orders.  This is only a formal compatibility test; it neither
constructs an arithmetic zero nor proves that further analytic
arguments cannot exclude it.

\paragraph{Sources, input class, and motive interpretation.}
The modular holomorphy input is Calegari (2013), Theorem~1.1,
publisher-PDF p.~2, and the exact all-prime induction identities
are the previously checked MT-S506--MT-S508 locators above.
The classical completed Hecke continuation and functional
equation are Ramakrishnan and Valenza (1999), \S7.4,
Theorem~7-19, printed p.~279, and \S7.7, Theorem~7-28,
printed p.~289, as recorded in MT-M18; NT-Z10 supplies the
fixed-strip growth and zero count.  The necessary Hadamard
normalization, archimedean estimate, and optimization are
proved in full here.  No new zero-free-region theorem is an
input, and no priority claim is made for the classical
logarithmic-derivative method.  The input class of MA-RB-001
is $\mathsf A$, not $M$.

The regular-character identity in (MA-RB.6) is the
$L$-function identity of the regular Artin motive attached
to the finite \'etale cover $\operatorname{Spec}L\to
\operatorname{Spec}\mathbb Q$.  Its decomposition with
multiplicities equal to irreducible dimensions explains
the coefficient five in (MA-RB.7).
The exact available motive locators remain Deligne (1979),
Definition~6.1 and \S\S6.2--6.3, and SGA~1,
Expos\'e~V, \S7, already checked in MA-S5.
This interpretation supplies the factorization, not the
analytic inequality or a number-field purity assertion.

Heilbronn, H. (1973). On real zeros of Dedekind $\zeta$-functions.
\emph{Canadian Journal of Mathematics, 25}(4), 870--873.
\url{https://doi.org/10.4153/CJM-1973-090-3}.
